\chapter{Probabilities and the Interpretation of Conditionals}
\label{sec:p19-rationality-4-6}

Oberauer and Pessach's 2021 chapter asks two different questions: what an
ordinary conditional means, and how people reason from conditional premises.
Its case for a probabilistic interpretation is not a claim that one Bayesian
updating model explains every conditional inference. The chapter ultimately
favors a distinction between knowledge-sensitive and form-sensitive
contributions to inference, while retaining objections concerning relevance,
task wording and the status of conditional probabilities
\cite[pp.~295--304]{OberauerPessach2021ProbabilitiesConditionals}.

This treatment keeps formal standards, reports of psychological tasks and
the authors' normative evaluations separate. Its experimental comparisons
are the chapter's dated source reports unless an inspected original-study
record is expressly identified. Approximate review rates are not estimates
of how all people interpret every everyday conditional.
\Cref{sec:p19-evidence-methods} supplies the corresponding design,
denominator and generalization distinctions.

\section{What the conditional task asks}
\label{sec:p19-rat46-tasks}

Conditionals can describe causal, logical, normative or other dependencies,
and can express hypothetical deliberation. The source's opening examples
include considering the consequences of reducing emissions and asking what
follows if one chooses a course of action. Formalizing their language and
evaluating their empirical truth are additional questions. Here \(A\)
names the antecedent and \(C\) the consequent. The source writes
\(A\to C\) for the conditional under discussion. Within a probability
judgment this is its ordinary-language conditional, not automatically the
material-implication event. The material comparison is explicitly written
\(\neg A\lor C\) below.

\begin{longtable}{@{}p{.21\textwidth}p{.35\textwidth}p{.36\textwidth}@{}}
\caption{Distinct elicitation questions in the source treatment.}
\label{tab:p19-rat46-tasks}\\
\toprule
Task & Supplied information and requested answer & What the response does not settle\\
\midrule\endfirsthead
\caption[]{Distinct elicitation questions in the source treatment (continued)}\\
\toprule Task & Supplied information and requested answer & What the response does not settle\\
\midrule\endhead
\midrule\multicolumn{3}{r}{\emph{continued on next page}}\\
\endfoot
\endlastfoot
Truth-table judgment & An antecedent/consequent truth combination;
judge the conditional true, false or irrelevant. & Irrelevance is not
automatically bivalent falsity or a numerical probability.\\\addlinespace
Inference endorsement & A conditional, a categorical premise and a
proposed conclusion; judge whether it follows. & An endorsement frequency
is not a conditional probability defined from a joint distribution.\\\addlinespace
Probability judgment & Frequencies or background knowledge; estimate
how likely the conditional is. & A rating can reflect interpretation,
arithmetic, response format or other task demands.\\\addlinespace
Conditional-probability question & Supposing \(A\), assess \(C\). &
This is not automatically the probability of a fixed proposition
expressed by the entire conditional.\\\addlinespace
Acceptability or relevance & Assess whether a conditional is a
reasonable assertion in its context or has an inferential link. &
Assertion quality, conditional probability and truth are distinct measures.\\
\bottomrule
\end{longtable}

The classical truth table makes the material conditional false only in
the \(A\land\neg C\) case. It is true in the other three cases.
The chapter's Table 4.6.1 contrasts this with the majority truth-table
responses attributed to Johnson-Laird and Tagart (1969): true for
\(A\land C\), false for \(A\land\neg C\), and irrelevant for
both false-antecedent cases. Evans et al. (2008) and Schroyens (2010)
are additional cited sources for this response pattern. The chapter
does not supply those studies' recruitment, sample denominators or
uncertainty intervals; the comparison is not a new pooled study
\cite[pp.~295--296]{OberauerPessach2021ProbabilitiesConditionals}.

The material reading also licenses an arbitrary antecedent when the
consequent is true, and an arbitrary consequent when the antecedent is
false. The source illustrates the first with a conditional linking cats
being fish to cats being mammals, and the second with a square moon
antecedent and opposed evaluations of a president. These are contrasts
between a truth function and ordinary assertion, not new evidence about
cats, astronomy or political evaluation.

In the four familiar inference tasks, MP supplies \(A\) and proposes
\(C\); MT supplies \(\neg C\) and proposes \(\neg A\); AC supplies
\(C\) and proposes \(A\); DA supplies \(\neg A\) and proposes
\(\neg C\). The conditional is the common major premise. With fixed
material meanings, MP and MT are valid and AC/DA are not, as proved in
\cref{thm:p19-conditional-benchmark}. The 2021 review reports approximate
endorsement of 95\% for MP, 70\% for MT and more than 50\% for AC/DA,
citing the 1993 reasoning literature and Schroyens--Schaeken (2003).
These are broad literature summaries, not denominators from one experiment.
The authors use the mismatch, everyday uncertainty and the material
paradoxes to motivate a probabilistic alternative; they do not thereby
invalidate the classical theorems on their declared meanings.

\section{The Equation, supposition, and the truth question}
\label{sec:p19-rat46-equation}

\begin{hbdefinition}{def:p19-rat46-equation}{The source's probability-conditional assessment}
Oberauer and Pessach call
\[
                         P(A\to C)=P(C\mid A)
\]
the Equation. The left side denotes a degree-of-belief assessment of
an indicative conditional in the account being examined. It is not
defined as \(P(\neg A\lor C)\). The right side can denote either
the ordinary positive-antecedent ratio
\[
                  P(C\mid A)=\frac{P(A\land C)}{P(A)}
                  \qquad(P(A)>0),
\]
or the result of a suppositional assessment: temporarily suppose \(A\)
and assess \(C\) using the relevant knowledge. Which interpretation
is intended must be stated. Numerical equality is not an identification
of a conditional sentence with an event in the original Boolean algebra.
\end{hbdefinition}

The Ramsey-test motivation is to add the antecedent hypothetically and
consider the consequent from that hypothetical information state. The
chapter's source is Ramsey's 1929 text as reprinted in 1990; it invokes
degrees of belief rather than treating every inquiry as a categorical
true/false judgment. Its probabilities are subjective degrees of belief,
not necessarily objective frequencies of a population
\cite[pp.~296--297]{OberauerPessach2021ProbabilitiesConditionals}.

The authors distinguish assessment from explicit ratio computation.
Edgington's deliberation case considers whether to choose action \(A\)
or another action and what would follow, without first fixing a probability
that one will choose \(A\). It is possible to consider consequences of
a supposition without having a complete numerical unconditional joint
distribution. This does not change the ratio theorem for a coherent,
already specified probability space; it changes the information and
procedure used by the deliberator.

When \(P(A)=0\), the ordinary ratio is undefined. The chapter contrasts
Stalnaker's extension of the Ramsey test, which first adjusts beliefs to
make the antecedent compatible and then adds it, with Bennett's refusal
to assign a conditional probability in most such cases. Minimal adjustment
is itself contested and is not uniquely specified by the phrase
``suppose \(A\).'' A counterfactual antecedent and a zero-probability
event must not be silently identified in every probability model.
\Cref{def:p08-stalnaker-conditional} supplies a particular selection
semantics, and \cref{sec:p15-rationality-4-4} supplies a different,
primitive conditional-event/coherence treatment of null strata. Neither
is automatically the unique formalization of this chapter's psychological
Ramsey-test account.

\begin{hbtheorem}{thm:p19-rat46-fixed-event}{proved=annex:p19.rat46.fixed-event}{The finite fixed-event obstruction}
Let \(A,C\) generate a four-atom Boolean algebra with all four atoms
nonempty. No fixed event \(D\) in that algebra satisfies
\(P(D)=P(C\mid A)\) for every strictly positive probability on it.
The fixed-event, same-algebra and probability-richness hypotheses are
essential to this finite statement. It is a complete local illustration
of the chapter's truth-table argument, not a reconstruction of every
version of Lewis's triviality results.
\end{hbtheorem}

The chapter's broader discussion cites Lewis (1976) to reject an ordinary
proposition-plus-probability interpretation satisfying the Equation
uniformly. The finite result above explains why changing the truth-table
rows selected by a fixed Boolean proposition does not solve that problem.
It does not show that all conditional semantics are impossible, that a
three-valued conditional event is a bivalent Boolean event, or that an
assessment procedure cannot return the indicated number. The separately
attributed cited presentation of this finite boundary is
\cref{thm:p08-lewis-triviality}; no broader original Lewis theorem is
claimed to have been proved by this local illustration.

The authors' nonpropositional account still recognizes an antecedent-true
case as confirmed when \(C\) occurs and refuted when \(\neg C\)
occurs. With the antecedent not obtaining, its assessment is not the
ordinary bivalent truth of the material conditional. Asking how likely
the conditional is to be true \emph{given that it has an antecedent-true
case} differs from asking for the unconditional probability of
\(A\land C\). This is the distinction underlying their appeal to
Edgington and Jeffrey; it does not prove that every speaker gives that
interpretation to the words ``probability that it is true.''

\section{Examples and competing interpretations of the response}
\label{sec:p19-rat46-interpretations}

The probability conditional avoids treating a merely improbable antecedent
as a sufficient reason for high conditional belief: with \(P(A)>0\),
\(P(C\mid A)\) can still be small. Likewise high \(P(C)\) alone does
not force high \(P(C\mid A)\). The source contrasts Peter's usual
morning run with the unlikely claim that he runs when injured. If the
consequent is certain under the coherent probability model and the
antecedent is positive, however, its conditional probability is one.
The chapter notes that this latter boundary motivates a relevance objection.

Its London--France example stipulates incompatible locations for Peter.
With positive probability of London and no London-and-France cases,
\(P(C\mid A)=0\), regardless of whether London is Peter's actual
location. Calling the conditional false can then report extremely low
belief rather than assert that its antecedent actually occurred. The
chapter calls this a pleonastic use of truth/falsity, citing Jeffrey and
Over--Baratgin. Its note 3 applies the analogous observation to calling
a highly credible conditional true.

\begin{hbtheorem}{thm:p19-rat46-frequency-contrast}{proved=annex:p19.rat46.frequency-contrast}{The source's card and die contrasts}
For the chapter's illustrative 100 cards, with truth-case counts
\((A C,A\neg C,\neg A C,\neg A\neg C)=(50,10,20,20)\),
the conditional-probability, conjunction and material-conditional
values are respectively \(5/6,1/2,9/10\).
For its six die faces, two yellow circles, two yellow squares and two
red circles, the corresponding values for ``if yellow, circle'' are
\(1/2,1/3,2/3\). These are calculations under the stipulated uniform
draws, not participant-response measurements.
\end{hbtheorem}

The chapter contrasts its account with the 2002 basic-conditional version
of mental-model theory. A fully elaborated representation has the three
possible rows \(A C\), \(\neg A C\), \(\neg A\neg C\); a sparse
representation has the explicit \(A C\) case and unspecified alternatives.
The authors object that explaining false-antecedent irrelevance by failing
to elaborate those alternatives leaves unexplained how the same sparse
representation distinguishes the also-unrepresented falsifying
\(A\neg C\) case. This is their criticism of that version, not a
proof that no mental-model revision can distinguish the cases.
Their note 4 explicitly records revisions by Barrouillet et al. (2008)
and Johnson-Laird et al. (2015), and the associated criticisms by
Oberauer--Oaksford (2008) and Baratgin et al. (2015).
Johnson-Laird's separately selected 2021 treatment at
\cref{sec:p19-rationality-2-3} retains its own representation assumptions
and source date.

More direct tests compare conditional ratings with the ratio, conjunction
and material values. The chapter cites Evans--Handley--Over (2003)
and Oberauer--Wilhelm (2003) for abstract-card tasks; Over et al. (2007)
uses everyday conditionals and participants' own judgments of the four
conjunction probabilities, constrained to sum to 100\%. It reports strong
ratio correspondence and rejection of the material prediction, but also
a substantial conjunction-sensitive subgroup in abstract probability tasks.
These different elicitation procedures are not interchangeable measures
of a single, context-independent mental representation
\cite[pp.~298--299]{OberauerPessach2021ProbabilitiesConditionals}.

\paragraph{An original-study anchor for the probability versus truth contrast.}
Oberauer and Wilhelm's Experiments 1A/1B separated ratio and joint-frequency
effects in a \(2\times2\) design. Experiment 1A tested 61 Potsdam
school students aged 17--21 within participants; the Internet Experiment
1B accepted 2,255 initial responses and retained 2,002 after missing/illegal
ratings were removed. Its four retained cells were 539, 520, 468 and 475.
The information described 2,000 cards. Participants directly assessed
the conditional probability after an A-card was drawn, then estimated
the probability that a conditional held for a random sample of ten cards;
betting provided a further, utility-sensitive response format. Experiment
1B displayed the four conditions sequentially over about a year, rather
than randomly assigning all online respondents concurrently.

Ratio affected conditional ratings strongly, but joint frequency also
had a positive effect. The paper reports conditional ratings below the
direct conditional-probability ratings, despite their similar patterns.
Exact agreement between the two ratings occurred for 52\% of Experiment
1B participants; only 23\% of Experiment 1A participants matched them
consistently across all four conditions. Similar group patterns and exact
individual numerical agreement therefore have different denominators.
Its binary Experiment 2 lacked the corresponding joint-frequency effect
on the ordinary conditional, supporting the chapter's warning about task
format. The ten-card truth question is not literally a one-card event
query; the chapter's 100-card illustration above is not this experiment's
2,000-card material. No population prevalence or causal effect of web
recruitment follows from the reported results
\cite[pp.~682--686]{OberauerWilhelm2003MeaningsConditionals}.

\paragraph{What a conjunction pattern could mean.}
It does not establish that a person identifies every conditional with
conjunction. The source's wire-warning example shows why such a global
identification would undermine hypothetical consideration: not touching
the wire would make the conjunction false without refuting the warning.
Four explanations of the numerical pattern are considered:

\begin{itemize}
\item A sampled-subset task might encourage requiring at least one
      antecedent-and-consequent instance before calling the rule true.
\item A partial Ramsey-test computation might count \(A\land C\)
      but omit the opposing \(A\land\neg C\) cases. The opposing
      cases are not another copy of \(A\land C\).
\item Wording that asks for the probability the conditional is true may
      be read literally as counting its true cases, or instead as asking
      for its suppositional credibility. Edgington (2003) advances this
      ambiguity explanation; the chapter says it had not been directly tested.
\item Division and rescaling demands may favor the easier conjunction
      answer, even when another response format accesses a ratio-sensitive
      interpretation.
\end{itemize}

The chapter reports that Oberauer et al. (2007) rejected the first two
explanations in their tests. It challenges the wording-only account using
Oberauer--Wilhelm's binary-task result, which did not elicit the predicted
conjunction pattern. Everyday conditional ratings in Over et al. (2007)
and facilitated die responses in Fugard et al. (2011) favor an arithmetic-
load contribution. Its note 6 stresses that using background knowledge
need not mean explicitly computing two conjunction probabilities and a
ratio. These are rival explanations and source-reported constraints, not
independent demonstrations of one universally operating mechanism.

The Fugard study's published report has two Salzburg student experiments
(66 and 65 participants), 71 trials, and 46 uniquely classifiable items.
Modal conditional-event responses were 50/66 and 45/65; 36/65 participants
in Experiment 2 shifted to that response class, 29 from conjunction.
The chapter's approximate 73\% summary and its description of 70 trials
must not be combined with these distinct denominators. The facilitated
``out of'' format reduces division/rescaling demands; task adaptation,
classification choices and entity--feature order limit transfer from these
responses to everyday language. The source reports increasing conditional-
event responses, not that every participant has one permanent semantics
\cite[pp.~637--642,647]{FugardEtAl2011ConditionalEvent}.

\section{Relevance and the inferentialist challenge}
\label{sec:p19-rat46-relevance}

The source's missing-link conditional connects Peter wearing a blue shirt
with rising sea levels. If the shirt event is treated as probabilistically
independent of a highly likely sea-level event, the conditional probability
is high without supplying an informative connection. High consequent
probability \emph{alone} is not enough for this result; independence or
the actual joint assessments matter. This distinguishes the objection
from the preceding injured-run example.

Inferentialism demands a sufficiently strong argument from the antecedent
and relevant background to the consequent; the connection can be deductive,
inductive or abductive. The chapter cites Douven (2015), Krzy\.{z}anowska
(2015), and Krzy\.{z}anowska--Wenmackers--Douven (2014) for the theory.
It reports Skovgaard-Olsen et al. (2016) as finding that the Equation fit
regular conditionals much better than missing-link cases, with individual
variation; the 2017 scorekeeping paper's 17\% is its reported subgroup,
not a population estimate. The other 2017 relevance/reason-relations paper
also bears on whether centering survives a missing inferential link.

The corrected 2016 accepted manuscript explicitly distinguishes probability
and acceptability tasks. Its recruited Internet sample of 577 was reduced
to 348 after language, comprehension, timing and seriousness exclusions;
its indicative-probability subgroup had 94 participants. Positive, negative
and irrelevant relations and four prior conditions were varied within
participants, while indicative/concessive wording and probability/
acceptability evaluation formed four between-participant groups. Contexts
and conditions were counterbalanced, and mixed models allowed repeated
participant and scenario responses. This is not the same design or
denominator as the 2017 17\% scorekeeping report; no common numerical
effect is pooled from them
\cite[accepted manuscript pp.~16--20]{SkovgaardOlsenSingmannKlauer2016Relevance}.

From \(A\land C\), both the material and the stipulated positive-
antecedent probability standards support a conditional to \(C\).
The inferentialist requirement can nevertheless reject a missing-link
assertion. The chapter cites Cruz et al. (2015) for centering endorsement,
and the Skovgaard-Olsen papers for relevance-sensitive limits. Its remaining
dispute is whether relevance affects semantics or pragmatics; Douven et al.
(2018) and Krzy\.{z}anowska--Collins--Hahn (2017) are the separately cited
positions. Probability fit alone does not settle that dispute
\cite[pp.~299--300]{OberauerPessach2021ProbabilitiesConditionals}.

\section{Reasoning models, conditional premises, and form}
\label{sec:p19-rat46-reasoning-models}

The source first reconstructs the mental-model prediction for the four
inferences: an explicit \(A C\) model supports MP and AC; adding
\(\neg A C\) permits rejecting AC/DA, and adding \(\neg A\neg C\)
supports MT as well. This is a representation/process proposal. The
authors then assess a different, purely Bayesian model in which the
minor premise conditions a background joint distribution.

\begin{hbdefinition}{def:p19-rat46-joint}{The source's three probability parameters}
For a coherent joint distribution set \(a=P(A)\), \(b=P(C)\),
and \(e=1-P(C\mid A)\), with \(0<a<1\) for the ratio-based
description. Here \(e\) is the exceptions parameter, not a measured
endorsement rate. The four masses are
\[
\begin{aligned}
 P(A\land C)&=a(1-e),&P(A\land\neg C)&=ae,\\
 P(\neg A\land C)&=b-a(1-e),&P(\neg A\land\neg C)&=1-b-ae.
\end{aligned}
\]
Nonnegativity requires \(0\leq e\leq1\) and
\(a(1-e)\leq b\leq1-ae\). The source's model associates MP, MT,
AC and DA with \(P(C\mid A)\), \(P(\neg A\mid\neg C)\),
\(P(A\mid C)\) and \(P(\neg C\mid\neg A)\), respectively.
This associates probability predictions with task responses; it does
not define logical validity by how often an inference is endorsed.
\end{hbdefinition}

\begin{hbtheorem}{thm:p19-rat46-model-parameters}{proved=annex:p19.rat46.model-parameters}{The inference probabilities and their boundary}
Under \cref{def:p19-rat46-joint}, with each indicated denominator
positive, the four model quantities are
\[
 1-e,\qquad 1-\frac{ae}{1-b},\qquad
 \frac{a(1-e)}b,\qquad\frac{1-b-ae}{1-a}.
\]
Small \(e\) does not by itself force high MT probability:
\(a=0.5,e=0.01,b=0.994\) gives MP \(0.99\) but MT \(1/6\),
with all four masses positive. Bounds on the base-rate factor
\(a/(1-b)\), or an actual fitted-data restriction, are needed for
that stronger implication.
\end{hbtheorem}

Oberauer (2006) and Schroyens--Schaeken (2003) are cited for the original
model's empirical MP--MT fitting difficulty: fitted parameters underpredicted
MP and overpredicted MT. This is a reported model-comparison result,
not the unconditional theorem that any high MP probability forces high MT.
Oaksford--Chater's 2007/2013 revision uses a new exceptions parameter
after a counterexample update for MT/AC/DA while retaining the old one
for MP. The parameters are thereby decoupled. Oberauer and Pessach object
that a negative antecedent or negative consequent does not itself identify
an \(A\land\neg C\) counterexample: knowing a person was not hit
or was not injured does not alone refute the proposed hit--injury link.
They judge that justification less principled even if it improves fit.
That is their normative criticism of the update rationale, not a claim
that fitting or any belief revision is mathematically impossible.

The major conditional premise creates a further question: what distinguishes
inferring \(C\) from \(A\) alone from inferring it from both \(A\)
and the conditional? The original background distribution gives the same
\(P(C\mid A)\) unless the conditional changes the distribution or
the response process. Ordinary conditioning on a categorical event does
not automatically implement learning a probability constraint. The chapter
contrasts decreasing the exceptions parameter with Hartmann--Rafiee Rad's
Kullback--Leibler-minimizing adjustment of the joint distribution, and
mentions Bayesian-network approaches without identifying them with the
three-parameter model. The indicated network account is developed separately
in \cref{sec:p16-rationality-4-2}; the 2012 presentation is not treated
as a uniquely specified general updating theorem here.

Klauer--Beller--H\"utter (2010) contrast reduced inferences lacking the
major premise with full conditional arguments. The dual-source model uses
the probabilistic model for its background-knowledge contribution, four
form parameters for MP/MT/AC/DA judgments, and a weight for form relative
to knowledge. A wet-streets judgment can draw on a rain connection without
an explicitly presented conditional, while presenting one can make the
argument form salient. The chapter reports Singmann--Klauer--Beller (2016)
as finding that DSM explained more variance than the compared original/
revised Bayesian and network models and confirmed a distinctive prediction.
It also reports selective changes of form parameters under changes to
biconditionals, only-if wording or added premises, but stability when
form is held fixed. No raw sample, residual variance or fit uncertainty
for that comparison is supplied by the 2021 chapter, so its qualitative
review claim is not a reconstructed quantitative replication.

The authors emphasize that form sensitivity need not mean bivalent material
semantics. Probability logic supplies its own standards; those standards
are distinct from the four background conditional probabilities just listed.
The source invokes MP/MT versus AC/DA in this connection. The actual
coherence language, null-stratum policies and propagation results are in
\cref{sec:p15-rationality-4-4}, not supplied merely by an endorsement model.
DSM is a measurement model: it measures a form contribution without fully
explaining the process that produces it.

Dual-process theories offer a possible process interpretation: Type 1 is
fast and knowledge-sensitive, while Type 2 is more controlled and depends
on working memory. The chapter cites Verschueren et al. (2005) for a
probabilistic/mental-model combination and Oberauer (2006) for its model
comparison. It does not establish that the two DSM sources are literally
the two processes or that every form response requires slow processing.

Geiger--Oberauer's 2010 procedural proposal preserves hypothetical models
without making the conditional's meaning a full material truth table.
The conditional instructs the reasoner to add \(C\) to each model
containing \(A\). MP starts with a model of the categorical \(A\)
and obtains \(A C\). MT starts with \(\neg C\), tests augmentations
by \(A\) and by \(\neg A\), and rejects the former when the
conditional instruction adds conflicting \(C\); the surviving
\(\neg A\neg C\) model supports \(\neg A\). This is the
source's proposed process sketch, not a new experiment or an adequacy
proof for a registered calculus
\cite[pp.~300--302]{OberauerPessach2021ProbabilitiesConditionals}.

\section{What the source conclusion warrants}
\label{sec:p19-rat46-boundaries}

The chapter concludes that hypothetical consideration can produce a coarse
degree of confidence without explicit division: represent a supposition,
use relevant knowledge or simulation, and express the resulting uncertainty.
It cites Hegarty (2004) for mental simulation and Khemlani--Lotstein--
Johnson-Laird (2015) for coarse model-based estimates. This provides a
plausible interpretation of the Equation, not an observation that every
participant computed a ratio internally.

Oberauer and Pessach regard the Equation's support as stronger than the
support for purely Bayesian models of inference. Their normative conclusion
is that the latter lack a rationally justified way to update on a conditional
premise in the accounts examined; their descriptive conclusion favors
dual-source/process developments but leaves the form-based process in need
of further specification. These are their 2021 judgments, not a claim that
no later or richer Bayesian account can represent conditional information.
\Cref{sec:p19-rationality-4-5} preserves Chater--Oaksford's differently
attributed Bayesian proposal rather than merging the disagreement away.

The conditional-interpretation and suppression experiments in
\cref{sec:p19-conditional-interpretation-experiments,
sec:p19-conditional-reasoning-suppression} ask different questions from
these probability-of-conditional tasks. Extra-premise effects, binary
necessity judgments, numerical ratings and contextual acceptability need
their own denominators and formal benchmarks. Neither a good model fit
nor a response called an error decides the semantics without specifying
the task and its interpretation.

The review reports above retain their source extent and uncertainty limits.
Except for the explicitly inspected original-study anchors, the chapter
does not print the dependent studies' recruitment, complete exclusions,
participant-level response data or uncertainty estimates. Their approximate
rates and comparative judgments cannot responsibly be pooled into a new
population statistic or promoted into a universal normative prescription.
